\section{Conclusion}
\label{sec:conclusion}

We developed a provider-neutral runtime for large-scale agentic simulation that makes model proposals, repair, policy completion, and retained model autonomy explicit and measurable, together with a causal activation model, an atomic commit protocol, and a scheduling-strategy ladder ranging from sequential to capability- and queue-aware execution. We evaluated the runtime on two real HPC systems with different accelerator vendors, running an identical model revision on two multi-role coordination workloads under real concurrent inference load.

The strongest measured result is that retained model autonomy remains at 0.00--0.02\% across every system, workload, and configuration studied, even though 99.99\% of steps remain useful. The runtime's contract and policy-completion layer, not the underlying model, is responsible for that reliability, and this study reports the two figures side by side rather than let a single throughput number conceal the difference between them. A second result, obtained by treating an initial null finding as a question rather than an answer, is that two independently discovered measurement confounds, an insufficiently loaded evaluation harness and a hardcoded batching cap, can each invert a scheduling or serving-configuration comparison. Correcting both resolved a preregistered scheduling decision gate from an apparent large regression into a genuine statistical tie on both systems.

For practitioners evaluating LLM-serving or agent-scheduling configurations on real infrastructure, the practical implication is that a confidently non-overlapping result is not sufficient evidence of a real effect without confirming that the tested load actually stresses the parameter under test. Extending this evaluation to a workload with genuine multi-provider heterogeneity is the most direct next step toward determining whether capability-aware and queue-aware scheduling deliver a measurable benefit beyond the one workload class tested here.
